% Computing M(Otilde) for p=7
% O1 : 0
% a : 0
% b0 : 10
% gs : 80
% O2 : 20
% b : 10
% O3p, O3n : 60
% O3_cp0 : 40
% O3_cp : 0
% O3_cn0 : 50
% O3_cn : 0
% Qp : 10
% Qps, Qpi : 0
% Qn : 0
% Qns, Qni : 0
% Qps_theta, Qns_theta : 0
% Qps_spin, Qns_spin : 0
% Qmps, Qmns : 0
% delta : 0
% twistmat : 0
% TOTAL : 280
\section{Example: level $7^2$}
Let $\H=\H(-1,-7)$,
the quaternion algebra ramified precisely at $\infty$ and $7$.
A maximal order,
having a unique left ideal class, is given by
\begin{align*}
   \O=\id{a_{1}}&=\<1,i,\frac{1+j}{2},\frac{i+k}{2}>. \\
\end{align*}
Its index $p$ suborder is given by
\[
  \Ot = \<1,7i,\frac{1+j}{2},\frac{7i+k}{2}>;
\]
inequivalent $\Ot$-subideals for the $\O$-ideal are show in
Table~\ref{table:psi:7}.

We fix two index $p$ suborders of $\Ot$
\begin{align*}
  \Ops & = \<1,7i,\frac{1+j}{2},\frac{7i+7k}{2}>, \\
  \Ons & = \<1,7i,\frac{1+7j}{2},\frac{1+7i+5j+k}{2}>
\end{align*}
in the $+$ and $-$ genus respectively.
Table~\ref{table:psi:7} shows the maps from $\Ot$-ideals
to ternary quadratic forms of level $7^2$ in the $+$ genus
and in the $-$ genus, computed via $\Ops$- and $\Ons$-subideals respectively.
\begin{table}
\begin{tabular}{|c|l|c|c|c|}
\hline
$\O$-ideals & $\Ot$-subideals & $\chi$ & $+$ genus & $-$ genus \\
\hline\hline
$\id{a_{1}}$ & $\id{b_{1,1}}=\<1,7i,\frac{1+j}{2},\frac{7i+k}{2}>$ & $+$ & $Q_{1}^{+}$ & $Q_{1}^{-}$\\
 & $\id{b_{1,2}}=\<7,1+i,\frac{7+j}{2},\frac{8+i+k}{2}>$ & $+$ & $Q_{2}^{+}$ & $Q_{2}^{-}$\\
 & $\id{b_{1,3}}=\<7,3+i,\frac{7+j}{2},\frac{10+i+k}{2}>$ & $-$ & $Q_{3}^{+}$ & $Q_{2}^{-}$\\
 & $\id{b_{1,4}}=\<7,5+i,\frac{7+j}{2},\frac{12+i+k}{2}>$ & $-$ & $Q_{3}^{+}$ & $Q_{1}^{-}$\\
\hline
\end{tabular}
\caption{Maps $\Gross_{\Ops}$ and $\Gross_{\Ons}$ from the $\Ot$-Ideals
to ternary quadratic forms in the $+$ and $-$ genus respectively, level $7^2$.
          \label{table:psi:7}}
\end{table}
The actual coefficients of the ternary quadratic forms are given in
Table~\ref{table:qf3:7}, with the notation
as in~\eqref{eq:notation:qf3}.
\begin{table}
\begin{tabular}{|l|rrrrrr|}
\hline
 & $a_1$ & $a_2$ & $a_3$ & $a_{23}$ & $a_{13}$ & $a_{12}$ \\
\hline\hline
$Q_{1}^{+}$ & $1$,&$28$,&$56$,&$-28$,&$0$,&$0$ \\
$Q_{2}^{+}$ & $4$,&$8$,&$49$,&$0$,&$0$,&$-4$ \\
$Q_{3}^{+}$ & $8$,&$9$,&$25$,&$2$,&$4$,&$8$ \\
\hline
$Q_{1}^{-}$ & $12$,&$12$,&$13$,&$-8$,&$-8$,&$-4$ \\
$Q_{2}^{-}$ & $5$,&$17$,&$17$,&$6$,&$2$,&$2$ \\
\hline
\end{tabular}
\caption{Coefficients of ternary quadratic forms, level $7^2$.
          \label{table:qf3:7}}
\end{table}
% Computing M(Otilde)^f for f=49A
% sub : 80
% dps, dns : 0
% cps, cns : 0
% lps, lns : 25330
% E49A x 2  : [1, 0, 0, 0, 4, 0, 0, 0]
%   e_1 = [1, -1, 0, 0]~  --->  {+} 1  |   {+} 1
%   e_2 = [0, 0, 1, -1]~  --->  {+} 0  |   {+} -1
%   ker+    = [0; 0; 1; -1]
%   ef+     = [1, -1, 0, 0]~   (height = 2)
%   Theta+  = q-2q^{8}+q^{9}-2q^{16}+q^{25}+O(q^{29})
%   ker-    = [1; -1; 1; -1]
%   ef-     = [1/2, -1/2, -1/2, 1/2]~   (height = 1)
%   Theta-  = -q^{5}+2q^{12}+q^{13}+q^{17}-2q^{20}+O(q^{24})
\subsection{$\mf{49A}$}
By computing the Brandt matrices for $\Ot$,
we find the space $\M(\Ot)^{\mf{49A}}
$
of dimension $2$, spanned by
\begin{align*}
  \e{\mf{49A}}^+ = &\ \cls{b_{1,1}}-\cls{b_{1,2}},\\
\intertext{and}
  \e{\mf{49A}}^- = &\ \frac{\cls{b_{1,1}}-\cls{b_{1,2}}-\cls{b_{1,3}}+\cls{b_{1,4}}}{2},\\
\end{align*}
with heights $\,\<\e{\mf{49A}}^+,\e{\mf{49A}}^+>=2\,\<\e{\mf{49A}}^-,\e{\mf{49A}}^->=2$.
Using table~\ref{table:psi:7},
we see that
\begin{align*}
  \Gross^+_{\mf{49A}} & =  Q_{1}^{+} - Q_{2}^{+}, \\
\intertext{and}
  \Gross^-_{\mf{49A}} & =  Q_{1}^{-} - Q_{2}^{-}. \\
\end{align*}
\begin{table}
\begin{tabular}{||r|rr||r|rr||}
\hline
$d$ & $c^{}_{\mf{49A}}(d)$ & $L(\mf{49A},d,1)$& $d$ & $c^{}_{\mf{49A}}(d)$ & $L(\mf{49A},d,1)$\\
\hline
   1 & 1 & 0.966656 & 109 & 2 & 0.370355 \\
   8 & -2 & 1.367058 & 113 & -4 & 1.454965 \\
   29 & 2 & 0.718014 & 120 & 4 & 1.411891 \\
   37 & 2 & 0.635669 & 137 & 2 & 0.330348 \\
   44 & 0 & 0.000000 & 141 & -4 & 1.302514 \\
   53 & 0 & 0.000000 & 149 & 0 & 0.000000 \\
   57 & 0 & 0.000000 & 156 & 4 & 1.238311 \\
   60 & 4 & 1.996716 & 165 & -4 & 1.204065 \\
   65 & -2 & 0.479596 & 172 & 0 & 0.000000 \\
   85 & -2 & 0.419394 & 177 & -8 & 4.650131 \\
   88 & -4 & 1.648734 & 184 & -4 & 1.140205 \\
   92 & 4 & 1.612493 & 193 & 0 & 0.000000 \\
   93 & -4 & 1.603801 & 197 & 0 & 0.000000 \\
\hline
$d$ & $c^{}_{\mf{49A}}(d)$ & $L(\mf{49A},d,1)$& $d$ & $c^{}_{\mf{49A}}(d)$ & $L(\mf{49A},d,1)$\\
\hline
   5 & -1 & 0.864603 & 97 & -1 & 0.196298 \\
   12 & 2 & 2.232396 & 101 & -1 & 0.192372 \\
   13 & 1 & 0.536204 & 104 & 4 & 3.033229 \\
   17 & 1 & 0.468897 & 124 & -4 & 2.777864 \\
   24 & -2 & 1.578542 & 129 & 4 & 2.723498 \\
   33 & -2 & 1.346185 & 136 & 2 & 0.663120 \\
   40 & 0 & 0.000000 & 145 & 2 & 0.642211 \\
   41 & 1 & 0.301933 & 152 & -2 & 0.627249 \\
   61 & 1 & 0.247535 & 157 & -3 & 1.388656 \\
   69 & -2 & 0.930974 & 173 & 1 & 0.146987 \\
   73 & -3 & 2.036493 & 181 & 1 & 0.143702 \\
   76 & -2 & 0.887064 & 185 & -4 & 2.274238 \\
   89 & 3 & 1.844376 & 188 & 0 & 0.000000 \\
\hline
\end{tabular}
\caption{Coefficients of
$\Gross^+_{\mf{49A}}$ (top),
$\Gross^-_{\mf{49A}}$ (bottom), and central values for $\mf{49A}$
          \label{table:dlcs:49A}}
\end{table}
Table~\ref{table:dlcs:49A} shows
the values of $c^{}_{\mf{49A}}(d)$
and $L(\mf{49A},-7d,1) = L(\mf{49A},d,1)$,
where $0<d<200$ is a fundamental discriminant
such that
$7\nmid d$.
The formula
\[
  L(\mf{49A},-7d,1) = L(\mf{49A},d,1)
  = k_{\mf{49A}} \frac{c_{\mf{49A}}(d)^2}{\sqrt{d}}
\cdot\begin{cases}
      1 & \text{ if $\kro{d}{7} = +1$} \\
      2 & \text{ if $\kro{d}{7} = -1$} \\
    \end{cases}
\]
is satisfied, where 
\[
  k_{\mf{49A}} = \frac{1}{4}\cdot
  \frac{\<\mf{49A},\mf{49A}>}
        {L(\mf{49A},1)\sqrt{7}}
  = 0.9666558528084057733665384189  = L(\mf{49A},1).
\]
% Computing M(Otilde) for p=11
% O1 : 0
% a : 10
% b0 : 20
% gs : 240
% O2 : 30
% b : 20
% O3p, O3n : 90
% O3_cp0 : 60
% O3_cp : 10
% O3_cn0 : 70
% O3_cn : 20
% Qp : 10
% Qps, Qpi : 0
% Qn : 10
% Qns, Qni : 0
% Qps_theta, Qns_theta : 0
% Qps_spin, Qns_spin : 10
% Qmps, Qmns : 0
% delta : 0
% twistmat : 0
% TOTAL : 600
\section{Example: level $11^2$}
Let $\H=\H(-1,-11)$,
the quaternion algebra ramified precisely at $\infty$ and $11$.
A maximal order,
and representatives for its left ideals classes, are given by
\begin{align*}
   \O=\id{a_{1}}&=\<1,i,\frac{1+j}{2},\frac{i+k}{2}>, \\
     \id{a_{2}}&=\<2,2i,\frac{3+2i+j}{2},\frac{2+3i+k}{2}>. \\
\end{align*}
Its index $p$ suborder is given by
\[
  \Ot = \<1,11i,\frac{1+j}{2},\frac{11i+k}{2}>;
\]
inequivalent $\Ot$-subideals for each $\O$-ideal are show in
Table~\ref{table:psi:11}.

We fix two index $p$ suborders of $\Ot$
\begin{align*}
  \Ops & = \<1,11i,\frac{1+j}{2},\frac{11i+11k}{2}>, \\
  \Ons & = \<1,11i,\frac{1+11j}{2},\frac{1+11i+j+k}{2}>
\end{align*}
in the $+$ and $-$ genus respectively.
Table~\ref{table:psi:11} shows the maps from $\Ot$-ideals
to ternary quadratic forms of level $11^2$ in the $+$ genus
and in the $-$ genus, computed via $\Ops$- and $\Ons$-subideals respectively.
\begin{table}
\begin{tabular}{|c|l|c|c|c|}
\hline
$\O$-ideals & $\Ot$-subideals & $\chi$ & $+$ genus & $-$ genus \\
\hline\hline
$\id{a_{1}}$ & $\id{b_{1,1}}=\<1,11i,\frac{1+j}{2},\frac{11i+k}{2}>$ & $+$ & $Q_{1}^{+}$ & $Q_{1}^{-}$\\
 & $\id{b_{1,2}}=\<11,5+i,\frac{11+j}{2},\frac{16+i+k}{2}>$ & $+$ & $Q_{2}^{+}$ & $Q_{2}^{-}$\\
 & $\id{b_{1,3}}=\<11,9+i,\frac{11+j}{2},\frac{20+i+k}{2}>$ & $+$ & $Q_{2}^{+}$ & $Q_{3}^{-}$\\
 & $\id{b_{1,4}}=\<11,4+i,\frac{11+j}{2},\frac{4+i+k}{2}>$ & $-$ & $Q_{3}^{+}$ & $Q_{3}^{-}$\\
 & $\id{b_{1,5}}=\<11,10+i,\frac{11+j}{2},\frac{10+i+k}{2}>$ & $-$ & $Q_{4}^{+}$ & $Q_{1}^{-}$\\
 & $\id{b_{1,6}}=\<11,3+i,\frac{11+j}{2},\frac{14+i+k}{2}>$ & $-$ & $Q_{3}^{+}$ & $Q_{2}^{-}$\\
\hline
$\id{a_{2}}$ & $\id{b_{2,1}}=\<22,14+2i,\frac{3+2i+j}{2},\frac{10+3i+k}{2}>$ & $+$ & $Q_{5}^{+}$ & $Q_{4}^{-}$\\
 & $\id{b_{2,2}}=\<22,8+2i,\frac{19+2i+j}{2},\frac{34+3i+k}{2}>$ & $+$ & $Q_{5}^{+}$ & $Q_{4}^{-}$\\
 & $\id{b_{2,3}}=\<22,12+2i,\frac{23+2i+j}{2},\frac{18+3i+k}{2}>$ & $-$ & $Q_{6}^{+}$ & $Q_{4}^{-}$\\
 & $\id{b_{2,4}}=\<2,22i,\frac{3+22i+j}{2},\frac{2+11i+k}{2}>$ & $-$ & $Q_{6}^{+}$ & $Q_{4}^{-}$\\
\hline
\end{tabular}
\caption{Maps $\Gross_{\Ops}$ and $\Gross_{\Ons}$ from the $\Ot$-Ideals
to ternary quadratic forms in the $+$ and $-$ genus respectively, level $11^2$.
          \label{table:psi:11}}
\end{table}
The actual coefficients of the ternary quadratic forms are given in
Table~\ref{table:qf3:11}, with the notation
as in~\eqref{eq:notation:qf3}.
\begin{table}
\begin{tabular}{|l|rrrrrr|}
\hline
 & $a_1$ & $a_2$ & $a_3$ & $a_{23}$ & $a_{13}$ & $a_{12}$ \\
\hline\hline
$Q_{1}^{+}$ & $1$,&$44$,&$132$,&$-44$,&$0$,&$0$ \\
$Q_{2}^{+}$ & $16$,&$16$,&$25$,&$-4$,&$-4$,&$-12$ \\
$Q_{3}^{+}$ & $5$,&$36$,&$36$,&$28$,&$4$,&$4$ \\
$Q_{4}^{+}$ & $4$,&$12$,&$121$,&$0$,&$0$,&$-4$ \\
$Q_{5}^{+}$ & $5$,&$9$,&$124$,&$-8$,&$-4$,&$-2$ \\
$Q_{6}^{+}$ & $4$,&$33$,&$45$,&$-22$,&$-4$,&$0$ \\
\hline
$Q_{1}^{-}$ & $8$,&$13$,&$61$,&$2$,&$4$,&$8$ \\
$Q_{2}^{-}$ & $13$,&$21$,&$21$,&$-2$,&$-6$,&$-6$ \\
$Q_{3}^{-}$ & $17$,&$21$,&$21$,&$-2$,&$-14$,&$-14$ \\
$Q_{4}^{-}$ & $13$,&$21$,&$24$,&$-16$,&$-4$,&$-6$ \\
\hline
\end{tabular}
\caption{Coefficients of ternary quadratic forms, level $11^2$.
          \label{table:qf3:11}}
\end{table}
% Computing M(Otilde)^f for f=11A
% sub : 400
% dps, dns : 0
% cps, cns : 0
% lps, lns : 43840
% E11A x 1  : [-2, -1, 1, -2, 1, 4, -2, 0]
%   e_1 = [2, 2, 2, 2, 2, 2, -3, -3, -3, -3]~  --->  {+} 2  |   {+} 2
%   ef+     = [1, 1, 1, 1, 1, 1, -3/2, -3/2, -3/2, -3/2]~   (height = 15)
%   Theta+  = q-q^{4}-q^{5}-2q^{9}-q^{12}+O(q^{20})
%   ef-     = [1, 1, 1, 1, 1, 1, -3/2, -3/2, -3/2, -3/2]~   (height = 15)
%   Theta-  = 2q^{8}+2q^{17}+2q^{21}-2q^{24}-2q^{28}+O(q^{32})
\subsection{$\mf{11A}$}
By computing the Brandt matrices for $\O$,
we find the space $\M(\Ot)^{\mf{11A}}
=\psi_{\Ot}\left(\M(\O)^{\mf{11A}}\right)
$
of dimension $1$, spanned by
\begin{align*}
  \e{\mf{11A}}^+ = \e{\mf{11A}}^- = &\ \psi_{\Ot}\left(\frac{\cls{a_{1}}-\cls{a_{2}}}{2}\right),\\
\end{align*}
with heights $\,\<\e{\mf{11A}}^+,\e{\mf{11A}}^+>=\,\<\e{\mf{11A}}^-,\e{\mf{11A}}^->=15$.
Using table~\ref{table:psi:11},
we see that
\begin{align*}
  \Gross^+_{\mf{11A}} & =  Q_{1}^{+} + 2Q_{2}^{+} + 2Q_{3}^{+} + Q_{4}^{+} - 3Q_{5}^{+} - 3Q_{6}^{+}, \\
\intertext{and}
  \Gross^-_{\mf{11A}} & =  2Q_{1}^{-} + 2Q_{2}^{-} + 2Q_{3}^{-} - 6Q_{4}^{-}. \\
\end{align*}
\begin{table}
\begin{tabular}{||r|rr||r|rr||}
\hline
$d$ & $c^{}_{\mf{121D}}(d)$ & $L(\mf{121D},d,1)$& $d$ & $c^{}_{\mf{121D}}(d)$ & $L(\mf{121D},d,1)$\\
\hline
   1 & 1 & 1.759399 & 113 & -1 & 0.165510 \\
   5 & -1 & 0.786827 & 124 & 3 & 1.421988 \\
   12 & -1 & 0.507895 & 133 & 2 & 0.610237 \\
   37 & 1 & 0.289243 & 136 & -2 & 0.603469 \\
   53 & -2 & 0.966688 & 137 & -1 & 0.150316 \\
   56 & -2 & 0.940438 & 141 & 2 & 0.592673 \\
   60 & 3 & 2.044237 & 152 & -2 & 0.570824 \\
   69 & -1 & 0.211807 & 157 & 1 & 0.140415 \\
   89 & 3 & 1.678463 & 168 & -2 & 0.542962 \\
   92 & -1 & 0.183430 & 177 & -7 & 6.479982 \\
   93 & 1 & 0.182441 & 181 & 1 & 0.130775 \\
   97 & -3 & 1.607759 & 185 & 3 & 1.164182 \\
   104 & 2 & 0.690093 & 188 & 2 & 0.513269 \\
\hline
$d$ & $c^{}_{\mf{121D}}(d)$ & $L(\mf{121D},d,1)$& $d$ & $c^{}_{\mf{121D}}(d)$ & $L(\mf{121D},d,1)$\\
\hline
   8 & 2 & 2.488166 & 101 & 2 & 0.700267 \\
   13 & 0 & 0.000000 & 105 & -2 & 0.686799 \\
   17 & 2 & 1.706868 & 109 & -2 & 0.674079 \\
   21 & 2 & 1.535729 & 120 & -2 & 0.642442 \\
   24 & -2 & 1.436543 & 129 & 2 & 0.619626 \\
   28 & -2 & 1.329981 & 140 & -2 & 0.594785 \\
   29 & 0 & 0.000000 & 145 & -4 & 2.337762 \\
   40 & 2 & 1.112742 & 149 & 2 & 0.576542 \\
   41 & -4 & 4.396351 & 156 & -4 & 2.253835 \\
   57 & 4 & 3.728610 & 161 & -2 & 0.554640 \\
   61 & 0 & 0.000000 & 172 & 2 & 0.536612 \\
   65 & 4 & 3.491625 & 173 & 2 & 0.535059 \\
   73 & 0 & 0.000000 & 184 & -2 & 0.518818 \\
   76 & 0 & 0.000000 & 193 & 4 & 2.026309 \\
   85 & -2 & 0.763334 & 197 & -2 & 0.501408 \\
\hline
\end{tabular}
\caption{Coefficients of
$\Gross^+_{\mf{11A}}$ (top),
$\Gross^-_{\mf{11A}}$ (bottom), and central values for $\mf{121D}$
          \label{table:dlcs:11A}}
\end{table}
Table~\ref{table:dlcs:11A} shows
the values of $c^{}_{\mf{11A}}(d)$
and $L(\mf{11A},-11d,1) = L(\mf{121D},d,1)$,
where $0<d<200$ is a fundamental discriminant
such that
$11\nmid d$.
The formula
\[
  L(\mf{11A},-11d,1) = L(\mf{121D},d,1)
  = k_{\mf{11A}} \frac{c_{\mf{11A}}(d)^2}{\sqrt{d}}
\]
is satisfied, where 
\[
  k_{\mf{11A}} = \frac{2}{5}\cdot
  \frac{\<\mf{11A},\mf{11A}>}
        {L(\mf{11A},1)\sqrt{11}}
  = 1.759399038662040141251585974  = L(\mf{121D},1).
\]
% Computing M(Otilde)^f for f=121A
% sub : 10
% dps, dns : 0
% cps, cns : 0
% lps, lns : 44790
% E121A x 2  : [-1, 2, 1, 2, 0, -1, 5, -6]
%   e_1 = [1, -1, 0, 1, -1, 0, 0, 0, 0, 0]~  --->  {+} 1  |   {+} -1
%   e_2 = [1, 1, -2, -1, -1, 2, 0, 0, 0, 0]~  --->  {+} 1  |   {+} 3
%   ker+    = [0; 1; -1; -1; 0; 1; 0; 0; 0; 0]
%   ef+     = [1, -1/2, -1/2, 1/2, -1, 1/2, 0, 0, 0, 0]~   (height = 3)
%   Theta+  = q+q^{5}+q^{9}-2q^{12}-2q^{16}+O(q^{20})
%   ker-    = [2; -1; -1; 1; -2; 1; 0; 0; 0; 0]
%   ef-     = [0, 1/2, -1/2, -1/2, 0, 1/2, 0, 0, 0, 0]~   (height = 1)
%   Theta-  = q^{13}-q^{17}-2q^{24}+2q^{28}-q^{29}+O(q^{40})
\subsection{$\mf{121A}$}
By computing the Brandt matrices for $\Ot$,
we find the space $\M(\Ot)^{\mf{121A}}
$
of dimension $2$, spanned by
\begin{align*}
  \e{\mf{121A}}^+ = &\ \frac{2\cls{b_{1,1}}-\cls{b_{1,2}}-\cls{b_{1,3}}+\cls{b_{1,4}}-2\cls{b_{1,5}}+\cls{b_{1,6}}}{2},\\
\intertext{and}
  \e{\mf{121A}}^- = &\ \frac{\cls{b_{1,2}}-\cls{b_{1,3}}-\cls{b_{1,4}}+\cls{b_{1,6}}}{2},\\
\end{align*}
with heights $\,\<\e{\mf{121A}}^+,\e{\mf{121A}}^+>=3\,\<\e{\mf{121A}}^-,\e{\mf{121A}}^->=3$.
Using table~\ref{table:psi:11},
we see that
\begin{align*}
  \Gross^+_{\mf{121A}} & =  Q_{1}^{+} - Q_{2}^{+} + Q_{3}^{+} - Q_{4}^{+}, \\
\intertext{and}
  \Gross^-_{\mf{121A}} & =  Q_{2}^{-} - Q_{3}^{-}. \\
\end{align*}
\begin{table}
\begin{tabular}{||r|rr||r|rr||}
\hline
$d$ & $c^{}_{\mf{121C}}(d)$ & $L(\mf{121C},d,1)$& $d$ & $c^{}_{\mf{121C}}(d)$ & $L(\mf{121C},d,1)$\\
\hline
   1 & 1 & 1.666157 & 113 & 1 & 0.156739 \\
   5 & 1 & 0.745128 & 124 & -2 & 0.598501 \\
   12 & -2 & 1.923912 & 133 & -2 & 0.577897 \\
   37 & -1 & 0.273915 & 136 & 4 & 2.285948 \\
   53 & 1 & 0.228864 & 137 & -2 & 0.569398 \\
   56 & -2 & 0.890598 & 141 & 2 & 0.561263 \\
   60 & -2 & 0.860400 & 152 & 2 & 0.540573 \\
   69 & 0 & 0.000000 & 157 & 0 & 0.000000 \\
   89 & 1 & 0.176612 & 168 & 4 & 2.056749 \\
   92 & -2 & 0.694835 & 177 & 2 & 0.500944 \\
   93 & 2 & 0.691090 & 181 & -3 & 1.114600 \\
   97 & -1 & 0.169173 & 185 & 1 & 0.122498 \\
   104 & -2 & 0.653521 & 188 & 2 & 0.486068 \\
\hline
$d$ & $c^{}_{\mf{121C}}(d)$ & $L(\mf{121C},d,1)$& $d$ & $c^{}_{\mf{121C}}(d)$ & $L(\mf{121C},d,1)$\\
\hline
   8 & 0 & 0.000000 & 101 & 0 & 0.000000 \\
   13 & 1 & 1.386326 & 105 & 0 & 0.000000 \\
   17 & -1 & 1.212307 & 109 & 1 & 0.478767 \\
   21 & 0 & 0.000000 & 120 & 2 & 1.825183 \\
   24 & -2 & 4.081234 & 129 & -2 & 1.760363 \\
   28 & 2 & 3.778489 & 140 & -2 & 1.689792 \\
   29 & -1 & 0.928193 & 145 & -1 & 0.415100 \\
   40 & 2 & 3.161310 & 149 & 1 & 0.409491 \\
   41 & 1 & 0.780630 & 156 & 2 & 1.600792 \\
   57 & 2 & 2.648255 & 161 & 2 & 1.575739 \\
   61 & -2 & 2.559954 & 172 & 0 & 0.000000 \\
   65 & 1 & 0.619984 & 173 & 0 & 0.000000 \\
   73 & -2 & 2.340107 & 184 & 2 & 1.473969 \\
   76 & -2 & 2.293456 & 193 & 1 & 0.359798 \\
   85 & 1 & 0.542160 & 197 & -1 & 0.356126 \\
\hline
\end{tabular}
\caption{Coefficients of
$\Gross^+_{\mf{121A}}$ (top),
$\Gross^-_{\mf{121A}}$ (bottom), and central values for $\mf{121C}$
          \label{table:dlcs:121A}}
\end{table}
Table~\ref{table:dlcs:121A} shows
the values of $c^{}_{\mf{121A}}(d)$
and $L(\mf{121A},-11d,1) = L(\mf{121C},d,1)$,
where $0<d<200$ is a fundamental discriminant
such that
$11\nmid d$.
The formula
\[
  L(\mf{121A},-11d,1) = L(\mf{121C},d,1)
  = k_{\mf{121A}} \frac{c_{\mf{121A}}(d)^2}{\sqrt{d}}
\cdot\begin{cases}
      1 & \text{ if $\kro{d}{11} = +1$} \\
      3 & \text{ if $\kro{d}{11} = -1$} \\
    \end{cases}
\]
is satisfied, where 
\[
  k_{\mf{121A}} = \frac{1}{6}\cdot
  \frac{\<\mf{121A},\mf{121A}>}
        {L(\mf{121A},1)\sqrt{11}}
  = 1.666156920394216089937692029  = L(\mf{121C},1).
\]
% Computing M(Otilde)^f for f=121B
% sub : 10
% dps, dns : 0
% cps, cns : 0
% lps, lns : 0
% E121B x 2  : [0, -1, -3, 0, 0, 0, 0, 0]
%   e_1 = [0, 0, 0, 0, 0, 0, 1, -1, 0, 0]~  ---> 
%   e_2 = [0, 0, 0, 0, 0, 0, 0, 0, 1, -1]~  ---> 
\subsection{$\mf{121B}$}
By computing the Brandt matrices for $\Ot$,
we find the space $\M(\Ot)^{\mf{121B}}
$
of dimension $2$.
Using table~\ref{table:psi:11},
we can check that
\[
  \Gross_{\Ops}\left(\M(\Ot)^{\mf{121B}}\right)  = 
  \Gross_{\Ons}\left(\M(\Ot)^{\mf{121B}}\right)  = 0,
\]
which is expected since $L(\mf{121B},1)=0$.
% Computing M(Otilde)^f for f=121C
% sub : 0
% dps, dns : 0
% cps, cns : 0
% lps, lns : 42870
% E121C x 2  : [1, 2, 1, -2, 0, 1, -5, 6]
%   e_1 = [1, 1, -2, 1, 1, -2, 0, 0, 0, 0]~  --->  {+} 1  |   {+} 1
%   e_2 = [1, -1, 0, -1, 1, 0, 0, 0, 0, 0]~  --->  {+} 1  |   {+} 1
%   ker+    = [0; 1; -1; 1; 0; -1; 0; 0; 0; 0]
%   ef+     = [1, -1/2, -1/2, -1/2, 1, -1/2, 0, 0, 0, 0]~   (height = 3)
%   Theta+  = q+2q^{4}-q^{5}+q^{9}+2q^{12}+O(q^{36})
%   ker-    = [0; 1; -1; 1; 0; -1; 0; 0; 0; 0]
%   ef-     = [1, -1/2, -1/2, -1/2, 1, -1/2, 0, 0, 0, 0]~   (height = 3)
%   Theta-  = 2q^{8}+3q^{13}-q^{17}-4q^{21}-2q^{24}+O(q^{28})
\subsection{$\mf{121C}$}
We readily find the space $\M(\Ot)^{\mf{121C}}
=\twist\left(\M(\Ot)^{\mf{121A}}\right)$
of dimension $2$, spanned by
\begin{align*}
  \e{\mf{121C}}^+ = \e{\mf{121C}}^- = &\ \frac{2\cls{b_{1,1}}-\cls{b_{1,2}}-\cls{b_{1,3}}-\cls{b_{1,4}}+2\cls{b_{1,5}}-\cls{b_{1,6}}}{2},\\
\intertext{and}
  &\ \cls{b_{1,2}}-\cls{b_{1,3}}+\cls{b_{1,4}}-\cls{b_{1,6}},\\
\end{align*}
with heights $\,\<\e{\mf{121C}}^+,\e{\mf{121C}}^+>=\,\<\e{\mf{121C}}^-,\e{\mf{121C}}^->=3$.
Using table~\ref{table:psi:11},
we see that
\begin{align*}
  \Gross^+_{\mf{121C}} & =  Q_{1}^{+} - Q_{2}^{+} - Q_{3}^{+} + Q_{4}^{+}, \\
\intertext{and}
  \Gross^-_{\mf{121C}} & =  2Q_{1}^{-} - Q_{2}^{-} - Q_{3}^{-}. \\
\end{align*}
\begin{table}
\begin{tabular}{||r|rr||r|rr||}
\hline
$d$ & $c^{}_{\mf{121A}}(d)$ & $L(\mf{121A},d,1)$& $d$ & $c^{}_{\mf{121A}}(d)$ & $L(\mf{121A},d,1)$\\
\hline
   1 & 1 & 1.019795 & 113 & -1 & 0.095934 \\
   5 & -1 & 0.456066 & 124 & -6 & 3.296890 \\
   12 & 2 & 1.177558 & 133 & 2 & 0.353710 \\
   37 & -5 & 4.191331 & 136 & -2 & 0.349787 \\
   53 & 1 & 0.140080 & 137 & 2 & 0.348508 \\
   56 & -2 & 0.545103 & 141 & 2 & 0.343529 \\
   60 & 6 & 4.739578 & 152 & -2 & 0.330865 \\
   69 & -4 & 1.964302 & 157 & 4 & 1.302216 \\
   89 & -3 & 0.972882 & 168 & 4 & 1.258862 \\
   92 & 2 & 0.425284 & 177 & 2 & 0.306610 \\
   93 & -2 & 0.422991 & 181 & 1 & 0.075801 \\
   97 & -3 & 0.931900 & 185 & 3 & 0.674791 \\
   104 & -4 & 1.599986 & 188 & 2 & 0.297505 \\
\hline
$d$ & $c^{}_{\mf{121A}}(d)$ & $L(\mf{121A},d,1)$& $d$ & $c^{}_{\mf{121A}}(d)$ & $L(\mf{121A},d,1)$\\
\hline
   8 & 2 & 1.442208 & 101 & 2 & 0.405894 \\
   13 & 3 & 2.545562 & 105 & 4 & 1.592349 \\
   17 & -1 & 0.247337 & 109 & 1 & 0.097679 \\
   21 & -4 & 3.560600 & 120 & -2 & 0.372376 \\
   24 & -2 & 0.832659 & 129 & 2 & 0.359152 \\
   28 & -2 & 0.770892 & 140 & -2 & 0.344754 \\
   29 & 3 & 1.704340 & 145 & 5 & 2.117234 \\
   40 & -4 & 2.579900 & 149 & -1 & 0.083545 \\
   41 & -1 & 0.159265 & 156 & 2 & 0.326596 \\
   57 & -2 & 0.540301 & 161 & -2 & 0.321484 \\
   61 & 0 & 0.000000 & 172 & 8 & 4.976552 \\
   65 & 1 & 0.126490 & 173 & 2 & 0.310134 \\
   73 & 0 & 0.000000 & 184 & 10 & 7.518027 \\
   76 & 6 & 4.211226 & 193 & -5 & 1.835161 \\
   85 & -5 & 2.765307 & 197 & 1 & 0.072657 \\
\hline
\end{tabular}
\caption{Coefficients of
$\Gross^+_{\mf{121C}}$ (top),
$\Gross^-_{\mf{121C}}$ (bottom), and central values for $\mf{121A}$
          \label{table:dlcs:121C}}
\end{table}
Table~\ref{table:dlcs:121C} shows
the values of $c^{}_{\mf{121C}}(d)$
and $L(\mf{121C},-11d,1) = L(\mf{121A},d,1)$,
where $0<d<200$ is a fundamental discriminant
such that
$11\nmid d$.
The formula
\[
  L(\mf{121C},-11d,1) = L(\mf{121A},d,1)
  = k_{\mf{121C}} \frac{c_{\mf{121C}}(d)^2}{\sqrt{d}}
\]
is satisfied, where 
\[
  k_{\mf{121C}} = \frac{1}{6}\cdot
  \frac{\<\mf{121C},\mf{121C}>}
        {L(\mf{121C},1)\sqrt{11}}
  = 1.019794861782916556837117278  = L(\mf{121A},1).
\]
% Computing M(Otilde)^f for f=121D
% sub : 20
% dps, dns : 0
% cps, cns : 0
% lps, lns : 6510
% E121D x 1  : [2, -1, 1, 2, 0, -4, 2, 0]
%   e_1 = [2, 2, 2, -2, -2, -2, -3, -3, 3, 3]~  --->  {+} 2
%   ef+     = [1, 1, 1, -1, -1, -1, -3/2, -3/2, 3/2, 3/2]~   (height = 15)
%   Theta+  = q+3q^{4}-5q^{5}-2q^{9}-5q^{12}+O(q^{16})
\subsection{$\mf{121D}$}
We readily find the space $\M(\Ot)^{\mf{121D}}
=\twist\left(\M(\Ot)^{\mf{11A}}\right)$
of dimension $1$, spanned by
\begin{align*}
  \e{\mf{121D}}^+ = &\ \twist\circ\psi_{\Ot}\left(\frac{\cls{a_{1}}-\cls{a_{2}}}{2}\right),\\
\end{align*}
with height $\<\e{\mf{121D}}^+,\e{\mf{121D}}^+>=15$.
Using table~\ref{table:psi:11},
we see that
\begin{align*}
  \Gross^+_{\mf{121D}} & =  Q_{1}^{+} + 2Q_{2}^{+} - 2Q_{3}^{+} - Q_{4}^{+} - 3Q_{5}^{+} + 3Q_{6}^{+}; \\
\end{align*}
on the other hand
$\Gross_{\Ons}\left(\M(\Ot)^{\mf{121D}}\right)  = 0$,
which is expected since $\sign(\mf{11A})= +1$.
\begin{table}
\begin{tabular}{||r|rr||r|rr||}
\hline
$d$ & $c^{}_{\mf{11A}}(d)$ & $L(\mf{11A},d,1)$& $d$ & $c^{}_{\mf{11A}}(d)$ & $L(\mf{11A},d,1)$\\
\hline
   1 & 1 & 0.253842 & 113 & -5 & 0.596986 \\
   5 & -5 & 2.838038 & 124 & -5 & 0.569892 \\
   12 & -5 & 1.831946 & 133 & 10 & 2.201088 \\
   37 & 5 & 1.043284 & 136 & 10 & 2.176676 \\
   53 & 10 & 3.486786 & 137 & -5 & 0.542179 \\
   56 & 10 & 3.392105 & 141 & -10 & 2.137734 \\
   60 & -5 & 0.819271 & 152 & -10 & 2.058929 \\
   69 & 15 & 6.875768 & 157 & -15 & 4.558227 \\
   89 & -5 & 0.672680 & 168 & 10 & 1.958432 \\
   92 & -5 & 0.661621 & 177 & 5 & 0.476998 \\
   93 & 5 & 0.658054 & 181 & -15 & 4.245281 \\
   97 & 5 & 0.644343 & 185 & -5 & 0.466571 \\
   104 & 10 & 2.489124 & 188 & -10 & 1.851332 \\
\hline
\end{tabular}
\caption{Coefficients of $\Gross^+_{\mf{121D}}$ and central values for $\mf{11A}$
          \label{table:dlcs:121D+}}
\end{table}
Table~\ref{table:dlcs:121D+} shows
the values of $c^{}_{\mf{121D}}(d)$
and $L(\mf{121D},-11d,1) = L(\mf{11A},d,1)$,
where $0<d<200$ is a fundamental discriminant
such that
$\kro{d}{11}=+1$.
The formula
\[
  L(\mf{121D},-11d,1) = L(\mf{11A},d,1)
  = k_{\mf{121D}} \frac{c_{\mf{121D}}(d)^2}{\sqrt{d}}
\]
is satisfied, where 
\[
  k_{\mf{121D}} = \frac{11}{300}\cdot
  \frac{\<\mf{121D},\mf{121D}>}
        {L(\mf{121D},1)\sqrt{11}}
  = 0.2538418608559106843377589233  = L(\mf{11A},1).
\]
% Computing M(Otilde) for p=17
% O1 : 0
% a : 10
% b0 : 20
% gs : 310
% O2 : 30
% b : 70
% O3p, O3n : 120
% O3_cp0 : 70
% O3_cp : 40
% O3_cn0 : 60
% O3_cn : 40
% Qp : 40
% Qps, Qpi : 0
% Qn : 40
% Qns, Qni : 0
% Qps_theta, Qns_theta : 0
% Qps_spin, Qns_spin : 20
% Qmps, Qmns : 0
% delta : 10
% twistmat : 0
% TOTAL : 880
\section{Example: level $17^2$}
Let $\H=\H(-17,-3)$,
the quaternion algebra ramified precisely at $\infty$ and $17$.
A maximal order,
and representatives for its left ideals classes, are given by
\begin{align*}
   \O=\id{a_{1}}&=\<1,i,\frac{1+j}{2},\frac{3i+2j+k}{6}>, \\
     \id{a_{2}}&=\<2,2i,\frac{3+2i+j}{2},\frac{3i+2j+k}{6}>. \\
\end{align*}
Its index $p$ suborder is given by
\[
  \Ot = \<1,i,\frac{1+17j}{2},\frac{3+3i+17j+k}{6}>;
\]
inequivalent $\Ot$-subideals for each $\O$-ideal are show in
Table~\ref{table:psi:17}.

We fix two index $p$ suborders of $\Ot$
\begin{align*}
  \Ops & = \<1,17i,\frac{1+17j}{2},\frac{3+33i+17j+k}{6}>, \\
  \Ons & = \<1,17i,\frac{1+17j}{2},\frac{3+99i+17j+k}{6}>
\end{align*}
in the $+$ and $-$ genus respectively.
Table~\ref{table:psi:17} shows the maps from $\Ot$-ideals
to ternary quadratic forms of level $17^2$ in the $+$ genus
and in the $-$ genus, computed via $\Ops$- and $\Ons$-subideals respectively.
\begin{table}
\begin{tabular}{|c|l|c|c|c|}
\hline
$\O$-ideals & $\Ot$-subideals & $\chi$ & $+$ genus & $-$ genus \\
\hline\hline
$\id{a_{1}}$ & $\id{b_{1,1}}=\<17,i,\frac{1+j}{2},\frac{36+3i+2j+k}{6}>$ & $+$ & $Q_{1}^{+}$ & $Q_{1}^{-}$\\
 & $\id{b_{1,2}}=\<17,i,\frac{25+j}{2},\frac{84+3i+2j+k}{6}>$ & $+$ & $Q_{2}^{+}$ & $Q_{2}^{-}$\\
 & $\id{b_{1,3}}=\<17,i,\frac{21+j}{2},\frac{42+3i+2j+k}{6}>$ & $+$ & $Q_{2}^{+}$ & $Q_{2}^{-}$\\
 & $\id{b_{1,4}}=\<17,i,\frac{31+j}{2},\frac{96+3i+2j+k}{6}>$ & $-$ & $Q_{1}^{+}$ & $Q_{1}^{-}$\\
 & $\id{b_{1,5}}=\<17,i,\frac{23+j}{2},\frac{12+3i+2j+k}{6}>$ & $-$ & $Q_{2}^{+}$ & $Q_{2}^{-}$\\
 & $\id{b_{1,6}}=\<17,i,\frac{29+j}{2},\frac{24+3i+2j+k}{6}>$ & $-$ & $Q_{2}^{+}$ & $Q_{2}^{-}$\\
\hline
$\id{a_{2}}$ & $\id{b_{2,1}}=\<2,2i,\frac{3+2i+17j}{2},\frac{9+9i+17j+k}{6}>$ & $+$ & $Q_{3}^{+}$ & $Q_{3}^{-}$\\
 & $\id{b_{2,2}}=\<34,2i,\frac{55+2i+j}{2},\frac{144+3i+2j+k}{6}>$ & $+$ & $Q_{4}^{+}$ & $Q_{4}^{-}$\\
 & $\id{b_{2,3}}=\<34,2i,\frac{59+2i+j}{2},\frac{84+3i+2j+k}{6}>$ & $+$ & $Q_{5}^{+}$ & $Q_{4}^{-}$\\
 & $\id{b_{2,4}}=\<34,2i,\frac{35+2i+j}{2},\frac{36+3i+2j+k}{6}>$ & $+$ & $Q_{6}^{+}$ & $Q_{3}^{-}$\\
 & $\id{b_{2,5}}=\<34,2i,\frac{7+2i+j}{2},\frac{48+3i+2j+k}{6}>$ & $+$ & $Q_{7}^{+}$ & $Q_{5}^{-}$\\
 & $\id{b_{2,6}}=\<34,2i,\frac{27+2i+j}{2},\frac{156+3i+2j+k}{6}>$ & $+$ & $Q_{7}^{+}$ & $Q_{6}^{-}$\\
 & $\id{b_{2,7}}=\<34,2i,\frac{67+2i+j}{2},\frac{168+3i+2j+k}{6}>$ & $+$ & $Q_{6}^{+}$ & $Q_{7}^{-}$\\
 & $\id{b_{2,8}}=\<34,2i,\frac{43+2i+j}{2},\frac{120+3i+2j+k}{6}>$ & $+$ & $Q_{5}^{+}$ & $Q_{6}^{-}$\\
 & $\id{b_{2,9}}=\<34,2i,\frac{47+2i+j}{2},\frac{60+3i+2j+k}{6}>$ & $+$ & $Q_{4}^{+}$ & $Q_{5}^{-}$\\
 & $\id{b_{2,10}}=\<34,2i,\frac{51+2i+j}{2},\frac{3i+2j+k}{6}>$ & $-$ & $Q_{3}^{+}$ & $Q_{3}^{-}$\\
 & $\id{b_{2,11}}=\<34,2i,\frac{63+2i+j}{2},\frac{24+3i+2j+k}{6}>$ & $-$ & $Q_{4}^{+}$ & $Q_{4}^{-}$\\
 & $\id{b_{2,12}}=\<34,2i,\frac{23+2i+j}{2},\frac{12+3i+2j+k}{6}>$ & $-$ & $Q_{5}^{+}$ & $Q_{4}^{-}$\\
 & $\id{b_{2,13}}=\<34,2i,\frac{31+2i+j}{2},\frac{96+3i+2j+k}{6}>$ & $-$ & $Q_{6}^{+}$ & $Q_{3}^{-}$\\
 & $\id{b_{2,14}}=\<34,2i,\frac{19+2i+j}{2},\frac{72+3i+2j+k}{6}>$ & $-$ & $Q_{7}^{+}$ & $Q_{5}^{-}$\\
 & $\id{b_{2,15}}=\<34,2i,\frac{15+2i+j}{2},\frac{132+3i+2j+k}{6}>$ & $-$ & $Q_{7}^{+}$ & $Q_{6}^{-}$\\
 & $\id{b_{2,16}}=\<34,2i,\frac{3+2i+j}{2},\frac{108+3i+2j+k}{6}>$ & $-$ & $Q_{6}^{+}$ & $Q_{7}^{-}$\\
 & $\id{b_{2,17}}=\<34,2i,\frac{11+2i+j}{2},\frac{192+3i+2j+k}{6}>$ & $-$ & $Q_{5}^{+}$ & $Q_{6}^{-}$\\
 & $\id{b_{2,18}}=\<34,2i,\frac{39+2i+j}{2},\frac{180+3i+2j+k}{6}>$ & $-$ & $Q_{4}^{+}$ & $Q_{5}^{-}$\\
\hline
\end{tabular}
\caption{Maps $\Gross_{\Ops}$ and $\Gross_{\Ons}$ from the $\Ot$-Ideals
to ternary quadratic forms in the $+$ and $-$ genus respectively, level $17^2$.
          \label{table:psi:17}}
\end{table}
The actual coefficients of the ternary quadratic forms are given in
Table~\ref{table:qf3:17}, with the notation
as in~\eqref{eq:notation:qf3}.
\begin{table}
\begin{tabular}{|l|rrrrrr|}
\hline
 & $a_1$ & $a_2$ & $a_3$ & $a_{23}$ & $a_{13}$ & $a_{12}$ \\
\hline\hline
$Q_{1}^{+}$ & $4$,&$51$,&$103$,&$-34$,&$-4$,&$0$ \\
$Q_{2}^{+}$ & $15$,&$32$,&$47$,&$24$,&$10$,&$4$ \\
$Q_{3}^{+}$ & $4$,&$52$,&$103$,&$36$,&$4$,&$4$ \\
$Q_{4}^{+}$ & $15$,&$35$,&$47$,&$-18$,&$-10$,&$-14$ \\
$Q_{5}^{+}$ & $8$,&$43$,&$60$,&$-16$,&$-4$,&$-4$ \\
$Q_{6}^{+}$ & $16$,&$32$,&$47$,&$-24$,&$-4$,&$-12$ \\
$Q_{7}^{+}$ & $15$,&$16$,&$100$,&$-12$,&$-4$,&$-12$ \\
\hline
$Q_{1}^{-}$ & $7$,&$11$,&$292$,&$-8$,&$-4$,&$-6$ \\
$Q_{2}^{-}$ & $23$,&$28$,&$40$,&$20$,&$12$,&$20$ \\
$Q_{3}^{-}$ & $7$,&$39$,&$79$,&$30$,&$6$,&$2$ \\
$Q_{4}^{-}$ & $12$,&$23$,&$75$,&$-10$,&$-8$,&$-4$ \\
$Q_{5}^{-}$ & $7$,&$20$,&$147$,&$-8$,&$-6$,&$-4$ \\
$Q_{6}^{-}$ & $27$,&$28$,&$39$,&$4$,&$26$,&$24$ \\
$Q_{7}^{-}$ & $3$,&$23$,&$295$,&$-22$,&$-2$,&$-2$ \\
\hline
\end{tabular}
\caption{Coefficients of ternary quadratic forms, level $17^2$.
          \label{table:qf3:17}}
\end{table}
% Computing M(Otilde)^f for f=17A
% sub : 2330
% dps, dns : 10
% cps, cns : 0
% lps, lns : 67230
% E17A x 1  : [-1, 0, -2, 4, 0, -2, 1, -4]
%   e_1 = [3, 3, 3, 3, 3, 3, -1, -1, -1, -1, -1, -1, -1, -1, -1, -1, -1, -1, -1, -1, -1, -1, -1, -1]~  --->  {+} 2  |   {+} 2
%   ef+     = [3/2, 3/2, 3/2, 3/2, 3/2, 3/2, -1/2, -1/2, -1/2, -1/2, -1/2, -1/2, -1/2, -1/2, -1/2, -1/2, -1/2, -1/2, -1/2, -1/2, -1/2, -1/2, -1/2, -1/2]~   (height = 18)
%   Theta+  = 2q^{4}-2q^{8}+2q^{15}-2q^{16}-2q^{19}+O(q^{32})
%   ef-     = [3/2, 3/2, 3/2, 3/2, 3/2, 3/2, -1/2, -1/2, -1/2, -1/2, -1/2, -1/2, -1/2, -1/2, -1/2, -1/2, -1/2, -1/2, -1/2, -1/2, -1/2, -1/2, -1/2, -1/2]~   (height = 18)
%   Theta-  = -q^{3}-q^{7}+3q^{11}-2q^{20}+q^{23}+O(q^{24})
\subsection{$\mf{17A}$}
By computing the Brandt matrices for $\O$,
we find the space $\M(\Ot)^{\mf{17A}}
=\psi_{\Ot}\left(\M(\O)^{\mf{17A}}\right)
$
of dimension $1$, spanned by
\begin{align*}
  \e{\mf{17A}}^+ = \e{\mf{17A}}^- = &\ \psi_{\Ot}\left(\frac{\cls{a_{1}}-\cls{a_{2}}}{2}\right),\\
\end{align*}
with heights $\,\<\e{\mf{17A}}^+,\e{\mf{17A}}^+>=\,\<\e{\mf{17A}}^-,\e{\mf{17A}}^->=18$.
Using table~\ref{table:psi:17},
we see that
\begin{align*}
  \Gross^+_{\mf{17A}} & =  3Q_{1}^{+} + 6Q_{2}^{+} - Q_{3}^{+} - 2Q_{4}^{+} - 2Q_{5}^{+} - 2Q_{6}^{+} - 2Q_{7}^{+}, \\
\intertext{and}
  \Gross^-_{\mf{17A}} & =  3Q_{1}^{-} + 6Q_{2}^{-} - 2Q_{3}^{-} - 2Q_{4}^{-} - 2Q_{5}^{-} - 2Q_{6}^{-} - Q_{7}^{-}. \\
\end{align*}
\begin{table}
\begin{tabular}{||r|rr||r|rr||}
\hline
$d$ & $c^{}_{\mf{289A}}(d)$ & $L(\mf{289A},-d,1)$& $d$ & $c^{}_{\mf{289A}}(d)$ & $L(\mf{289A},-d,1)$\\
\hline
   4 & 2 & 2.663758 & 104 & 4 & 2.089624 \\
   8 & -2 & 1.883561 & 111 & -2 & 0.505665 \\
   15 & 2 & 1.375559 & 115 & 2 & 0.496793 \\
   19 & -2 & 1.222216 & 120 & 0 & 0.000000 \\
   35 & -2 & 0.900515 & 123 & -6 & 4.323294 \\
   43 & 2 & 0.812439 & 127 & 2 & 0.472741 \\
   47 & 0 & 0.000000 & 132 & 0 & 0.000000 \\
   52 & 0 & 0.000000 & 151 & 2 & 0.433547 \\
   55 & 2 & 0.718362 & 152 & 0 & 0.000000 \\
   59 & -2 & 0.693584 & 155 & 2 & 0.427916 \\
   67 & 0 & 0.000000 & 168 & -4 & 1.644107 \\
   83 & -2 & 0.584771 & 179 & 2 & 0.398197 \\
   84 & 4 & 2.325119 & 183 & -2 & 0.393821 \\
   87 & 2 & 0.571170 & 191 & 0 & 0.000000 \\
   103 & 0 & 0.000000 & 195 & 4 & 1.526046 \\
\hline
$d$ & $c^{}_{\mf{289A}}(d)$ & $L(\mf{289A},-d,1)$& $d$ & $c^{}_{\mf{289A}}(d)$ & $L(\mf{289A},-d,1)$\\
\hline
   3 & -1 & 0.768961 & 95 & -4 & 2.186367 \\
   7 & -1 & 0.503403 & 107 & 3 & 1.158819 \\
   11 & 3 & 3.614190 & 116 & 2 & 0.494647 \\
   20 & -2 & 1.191269 & 131 & -3 & 1.047301 \\
   23 & 1 & 0.277716 & 139 & -1 & 0.112969 \\
   24 & 2 & 1.087475 & 143 & -2 & 0.445509 \\
   31 & -1 & 0.239213 & 148 & -2 & 0.437919 \\
   39 & -4 & 3.412341 & 159 & 2 & 0.422500 \\
   40 & 2 & 0.842354 & 163 & 7 & 5.111720 \\
   56 & 0 & 0.000000 & 164 & -4 & 1.664037 \\
   71 & 3 & 1.422585 & 167 & -1 & 0.103064 \\
   79 & -1 & 0.149848 & 184 & -8 & 6.283996 \\
   88 & -2 & 0.567915 & 199 & 1 & 0.094414 \\
   91 & -2 & 0.558475 & & &  \\
\hline
\end{tabular}
\caption{Coefficients of
$\Gross^+_{\mf{17A}}$ (top),
$\Gross^-_{\mf{17A}}$ (bottom), and central values for $\mf{289A}$
          \label{table:dlcs:17A}}
\end{table}
Table~\ref{table:dlcs:17A} shows
the values of $c^{}_{\mf{17A}}(d)$
and $L(\mf{17A},-17d,1) = L(\mf{289A},-d,1)$,
where $0>-d>-200$ is a fundamental discriminant
such that
$17\nmid d$.
The formula
\[
  L(\mf{17A},-17d,1) = L(\mf{289A},-d,1)
  = k_{\mf{17A}} \frac{c_{\mf{17A}}(d)^2}{\sqrt{d}}
\]
is satisfied, where 
\[
  k_{\mf{17A}} = \frac{1}{2}\cdot
  \frac{\<\mf{17A},\mf{17A}>}
        {L(\mf{17A},1)\sqrt{17}}
  = 1.331879106385216159220474762.
\]
% Computing M(Otilde)^f for f=289A
% sub : 0
% dps, dns : 10
% cps, cns : 0
% lps, lns : 0
% E289A x 1  : [-1, 0, 2, -4, 0, -2, 0, -4]
%   e_1 = [3, 3, 3, -3, -3, -3, -1, -1, -1, -1, -1, -1, -1, -1, -1, 1, 1, 1, 1, 1, 1, 1, 1, 1]~  ---> 
\subsection{$\mf{289A}$}
We readily find the space $\M(\Ot)^{\mf{289A}}
=\twist\left(\M(\Ot)^{\mf{17A}}\right)$
of dimension $1$.
Using table~\ref{table:psi:17},
we can check that
\[
  \Gross_{\Ops}\left(\M(\Ot)^{\mf{289A}}\right)  = 
  \Gross_{\Ons}\left(\M(\Ot)^{\mf{289A}}\right)  = 0,
\]
which is expected since $L(\mf{289A},1)=0$.
% Computing M(Otilde) for p=19
% O1 : 0
% a : 10
% b0 : 30
% gs : 310
% O2 : 20
% b : 100
% O3p, O3n : 120
% O3_cp0 : 70
% O3_cp : 50
% O3_cn0 : 80
% O3_cn : 50
% Qp : 40
% Qps, Qpi : 0
% Qn : 50
% Qns, Qni : 0
% Qps_theta, Qns_theta : 0
% Qps_spin, Qns_spin : 30
% Qmps, Qmns : 0
% delta : 10
% twistmat : 0
% TOTAL : 970
\section{Example: level $19^2$}
Let $\H=\H(-1,-19)$,
the quaternion algebra ramified precisely at $\infty$ and $19$.
A maximal order,
and representatives for its left ideals classes, are given by
\begin{align*}
   \O=\id{a_{1}}&=\<1,i,\frac{1+j}{2},\frac{i+k}{2}>, \\
     \id{a_{2}}&=\<2,2i,\frac{3+2i+j}{2},\frac{2+3i+k}{2}>. \\
\end{align*}
Its index $p$ suborder is given by
\[
  \Ot = \<1,19i,\frac{1+j}{2},\frac{19i+k}{2}>;
\]
inequivalent $\Ot$-subideals for each $\O$-ideal are show in
Table~\ref{table:psi:19}.

We fix two index $p$ suborders of $\Ot$
\begin{align*}
  \Ops & = \<1,19i,\frac{1+j}{2},\frac{19i+19k}{2}>, \\
  \Ons & = \<1,19i,\frac{1+19j}{2},\frac{19i+18j+k}{2}>
\end{align*}
in the $+$ and $-$ genus respectively.
Table~\ref{table:psi:19} shows the maps from $\Ot$-ideals
to ternary quadratic forms of level $19^2$ in the $+$ genus
and in the $-$ genus, computed via $\Ops$- and $\Ons$-subideals respectively.
\begin{table}
\begin{tabular}{|c|l|c|c|c|}
\hline
$\O$-ideals & $\Ot$-subideals & $\chi$ & $+$ genus & $-$ genus \\
\hline\hline
$\id{a_{1}}$ & $\id{b_{1,1}}=\<1,19i,\frac{1+j}{2},\frac{19i+k}{2}>$ & $+$ & $Q_{1}^{+}$ & $Q_{1}^{-}$\\
 & $\id{b_{1,2}}=\<19,14+i,\frac{19+j}{2},\frac{14+i+k}{2}>$ & $+$ & $Q_{2}^{+}$ & $Q_{2}^{-}$\\
 & $\id{b_{1,3}}=\<19,9+i,\frac{19+j}{2},\frac{28+i+k}{2}>$ & $+$ & $Q_{3}^{+}$ & $Q_{3}^{-}$\\
 & $\id{b_{1,4}}=\<19,17+i,\frac{19+j}{2},\frac{36+i+k}{2}>$ & $+$ & $Q_{3}^{+}$ & $Q_{4}^{-}$\\
 & $\id{b_{1,5}}=\<19,15+i,\frac{19+j}{2},\frac{34+i+k}{2}>$ & $+$ & $Q_{2}^{+}$ & $Q_{5}^{-}$\\
 & $\id{b_{1,6}}=\<19,16+i,\frac{19+j}{2},\frac{16+i+k}{2}>$ & $-$ & $Q_{4}^{+}$ & $Q_{4}^{-}$\\
 & $\id{b_{1,7}}=\<19,11+i,\frac{19+j}{2},\frac{30+i+k}{2}>$ & $-$ & $Q_{5}^{+}$ & $Q_{5}^{-}$\\
 & $\id{b_{1,8}}=\<19,1+i,\frac{19+j}{2},\frac{20+i+k}{2}>$ & $-$ & $Q_{6}^{+}$ & $Q_{1}^{-}$\\
 & $\id{b_{1,9}}=\<19,7+i,\frac{19+j}{2},\frac{26+i+k}{2}>$ & $-$ & $Q_{5}^{+}$ & $Q_{2}^{-}$\\
 & $\id{b_{1,10}}=\<19,6+i,\frac{19+j}{2},\frac{6+i+k}{2}>$ & $-$ & $Q_{4}^{+}$ & $Q_{3}^{-}$\\
\hline
$\id{a_{2}}$ & $\id{b_{2,1}}=\<38,2+2i,\frac{59+2i+j}{2},\frac{22+3i+k}{2}>$ & $+$ & $Q_{7}^{+}$ & $Q_{6}^{-}$\\
 & $\id{b_{2,2}}=\<38,14+2i,\frac{71+2i+j}{2},\frac{2+3i+k}{2}>$ & $+$ & $Q_{8}^{+}$ & $Q_{7}^{-}$\\
 & $\id{b_{2,3}}=\<38,12+2i,\frac{31+2i+j}{2},\frac{18+3i+k}{2}>$ & $+$ & $Q_{9}^{+}$ & $Q_{8}^{-}$\\
 & $\id{b_{2,4}}=\<38,26+2i,\frac{7+2i+j}{2},\frac{58+3i+k}{2}>$ & $+$ & $Q_{9}^{+}$ & $Q_{9}^{-}$\\
 & $\id{b_{2,5}}=\<38,24+2i,\frac{43+2i+j}{2},\frac{74+3i+k}{2}>$ & $+$ & $Q_{8}^{+}$ & $Q_{10}^{-}$\\
 & $\id{b_{2,6}}=\<38,36+2i,\frac{55+2i+j}{2},\frac{54+3i+k}{2}>$ & $+$ & $Q_{7}^{+}$ & $Q_{6}^{-}$\\
 & $\id{b_{2,7}}=\<38,16+2i,\frac{35+2i+j}{2},\frac{62+3i+k}{2}>$ & $+$ & $Q_{8}^{+}$ & $Q_{7}^{-}$\\
 & $\id{b_{2,8}}=\<38,6+2i,\frac{63+2i+j}{2},\frac{66+3i+k}{2}>$ & $+$ & $Q_{9}^{+}$ & $Q_{8}^{-}$\\
 & $\id{b_{2,9}}=\<38,32+2i,\frac{51+2i+j}{2},\frac{10+3i+k}{2}>$ & $+$ & $Q_{9}^{+}$ & $Q_{9}^{-}$\\
 & $\id{b_{2,10}}=\<38,22+2i,\frac{3+2i+j}{2},\frac{14+3i+k}{2}>$ & $+$ & $Q_{8}^{+}$ & $Q_{10}^{-}$\\
 & $\id{b_{2,11}}=\<38,18+2i,\frac{75+2i+j}{2},\frac{46+3i+k}{2}>$ & $-$ & $Q_{10}^{+}$ & $Q_{8}^{-}$\\
 & $\id{b_{2,12}}=\<38,34+2i,\frac{15+2i+j}{2},\frac{70+3i+k}{2}>$ & $-$ & $Q_{10}^{+}$ & $Q_{9}^{-}$\\
 & $\id{b_{2,13}}=\<38,30+2i,\frac{11+2i+j}{2},\frac{26+3i+k}{2}>$ & $-$ & $Q_{11}^{+}$ & $Q_{10}^{-}$\\
 & $\id{b_{2,14}}=\<38,2i,\frac{19+2i+j}{2},\frac{38+3i+k}{2}>$ & $-$ & $Q_{12}^{+}$ & $Q_{6}^{-}$\\
 & $\id{b_{2,15}}=\<38,8+2i,\frac{27+2i+j}{2},\frac{50+3i+k}{2}>$ & $-$ & $Q_{11}^{+}$ & $Q_{7}^{-}$\\
 & $\id{b_{2,16}}=\<38,4+2i,\frac{23+2i+j}{2},\frac{6+3i+k}{2}>$ & $-$ & $Q_{10}^{+}$ & $Q_{8}^{-}$\\
 & $\id{b_{2,17}}=\<38,20+2i,\frac{39+2i+j}{2},\frac{30+3i+k}{2}>$ & $-$ & $Q_{10}^{+}$ & $Q_{9}^{-}$\\
 & $\id{b_{2,18}}=\<38,10+2i,\frac{67+2i+j}{2},\frac{34+3i+k}{2}>$ & $-$ & $Q_{11}^{+}$ & $Q_{10}^{-}$\\
 & $\id{b_{2,19}}=\<2,38i,\frac{3+38i+j}{2},\frac{2+19i+k}{2}>$ & $-$ & $Q_{12}^{+}$ & $Q_{6}^{-}$\\
 & $\id{b_{2,20}}=\<38,28+2i,\frac{47+2i+j}{2},\frac{42+3i+k}{2}>$ & $-$ & $Q_{11}^{+}$ & $Q_{7}^{-}$\\
\hline
\end{tabular}
\caption{Maps $\Gross_{\Ops}$ and $\Gross_{\Ons}$ from the $\Ot$-Ideals
to ternary quadratic forms in the $+$ and $-$ genus respectively, level $19^2$.
          \label{table:psi:19}}
\end{table}
The actual coefficients of the ternary quadratic forms are given in
Table~\ref{table:qf3:19}, with the notation
as in~\eqref{eq:notation:qf3}.
\begin{table}
\begin{tabular}{|l|rrrrrr|}
\hline
 & $a_1$ & $a_2$ & $a_3$ & $a_{23}$ & $a_{13}$ & $a_{12}$ \\
\hline\hline
$Q_{1}^{+}$ & $1$,&$76$,&$380$,&$-76$,&$0$,&$0$ \\
$Q_{2}^{+}$ & $5$,&$76$,&$92$,&$-76$,&$-4$,&$0$ \\
$Q_{3}^{+}$ & $25$,&$36$,&$36$,&$-4$,&$-16$,&$-16$ \\
$Q_{4}^{+}$ & $17$,&$44$,&$44$,&$12$,&$16$,&$16$ \\
$Q_{5}^{+}$ & $20$,&$24$,&$73$,&$4$,&$8$,&$20$ \\
$Q_{6}^{+}$ & $4$,&$20$,&$361$,&$0$,&$0$,&$-4$ \\
$Q_{7}^{+}$ & $5$,&$16$,&$365$,&$16$,&$2$,&$4$ \\
$Q_{8}^{+}$ & $16$,&$24$,&$77$,&$20$,&$8$,&$4$ \\
$Q_{9}^{+}$ & $20$,&$36$,&$45$,&$20$,&$16$,&$12$ \\
$Q_{10}^{+}$ & $9$,&$44$,&$77$,&$28$,&$6$,&$8$ \\
$Q_{11}^{+}$ & $5$,&$61$,&$92$,&$16$,&$4$,&$2$ \\
$Q_{12}^{+}$ & $4$,&$77$,&$96$,&$40$,&$4$,&$4$ \\
\hline
$Q_{1}^{-}$ & $8$,&$21$,&$181$,&$2$,&$4$,&$8$ \\
$Q_{2}^{-}$ & $29$,&$29$,&$37$,&$-6$,&$-6$,&$-18$ \\
$Q_{3}^{-}$ & $13$,&$48$,&$48$,&$20$,&$8$,&$8$ \\
$Q_{4}^{-}$ & $29$,&$32$,&$32$,&$-12$,&$-8$,&$-8$ \\
$Q_{5}^{-}$ & $29$,&$29$,&$41$,&$-14$,&$-14$,&$-18$ \\
$Q_{6}^{-}$ & $21$,&$32$,&$53$,&$-20$,&$-14$,&$-16$ \\
$Q_{7}^{-}$ & $29$,&$32$,&$37$,&$-28$,&$-6$,&$-8$ \\
$Q_{8}^{-}$ & $12$,&$13$,&$184$,&$-12$,&$-4$,&$-4$ \\
$Q_{9}^{-}$ & $8$,&$29$,&$124$,&$20$,&$4$,&$4$ \\
$Q_{10}^{-}$ & $21$,&$29$,&$53$,&$26$,&$14$,&$2$ \\
\hline
\end{tabular}
\caption{Coefficients of ternary quadratic forms, level $19^2$.
          \label{table:qf3:19}}
\end{table}
% Computing M(Otilde)^f for f=19A
% sub : 3250
% dps, dns : 0
% cps, cns : 0
% lps, lns : 77200
% E19A x 1  : [0, -2, 3, -1, 3, -4, -3, 1]
%   e_1 = [2, 2, 2, 2, 2, 2, 2, 2, 2, 2, -1, -1, -1, -1, -1, -1, -1, -1, -1, -1, -1, -1, -1, -1, -1, -1, -1, -1, -1, -1]~  --->  {+} 2  |   {+} 2
%   ef+     = [1, 1, 1, 1, 1, 1, 1, 1, 1, 1, -1/2, -1/2, -1/2, -1/2, -1/2, -1/2, -1/2, -1/2, -1/2, -1/2, -1/2, -1/2, -1/2, -1/2, -1/2, -1/2, -1/2, -1/2, -1/2, -1/2]~   (height = 15)
%   Theta+  = q+q^{4}-q^{5}-q^{9}-2q^{16}+O(q^{17})
%   ef-     = [1, 1, 1, 1, 1, 1, 1, 1, 1, 1, -1/2, -1/2, -1/2, -1/2, -1/2, -1/2, -1/2, -1/2, -1/2, -1/2, -1/2, -1/2, -1/2, -1/2, -1/2, -1/2, -1/2, -1/2, -1/2, -1/2]~   (height = 15)
%   Theta-  = -2q^{12}-2q^{21}+2q^{29}-2q^{33}+2q^{37}+O(q^{40})
\subsection{$\mf{19A}$}
By computing the Brandt matrices for $\O$,
we find the space $\M(\Ot)^{\mf{19A}}
=\psi_{\Ot}\left(\M(\O)^{\mf{19A}}\right)
$
of dimension $1$, spanned by
\begin{align*}
  \e{\mf{19A}}^+ = \e{\mf{19A}}^- = &\ \psi_{\Ot}\left(\frac{\cls{a_{1}}-\cls{a_{2}}}{2}\right),\\
\end{align*}
with heights $\,\<\e{\mf{19A}}^+,\e{\mf{19A}}^+>=\,\<\e{\mf{19A}}^-,\e{\mf{19A}}^->=15$.
Using table~\ref{table:psi:19},
we see that
\begin{align*}
  \Gross^+_{\mf{19A}} & =  Q_{1}^{+} + 2Q_{2}^{+} + 2Q_{3}^{+} + 2Q_{4}^{+} + 2Q_{5}^{+} + Q_{6}^{+} \\ & - Q_{7}^{+} - 2Q_{8}^{+} - 2Q_{9}^{+} - 2Q_{10}^{+} - 2Q_{11}^{+} - Q_{12}^{+}, \\
\intertext{and}
  \Gross^-_{\mf{19A}} & =  2Q_{1}^{-} + 2Q_{2}^{-} + 2Q_{3}^{-} + 2Q_{4}^{-} + 2Q_{5}^{-} \\ & - 2Q_{6}^{-} - 2Q_{7}^{-} - 2Q_{8}^{-} - 2Q_{9}^{-} - 2Q_{10}^{-}. \\
\end{align*}
\begin{table}
\begin{tabular}{||r|rr||r|rr||}
\hline
$d$ & $c^{}_{\mf{361B}}(d)$ & $L(\mf{361B},d,1)$& $d$ & $c^{}_{\mf{361B}}(d)$ & $L(\mf{361B},d,1)$\\
\hline
   1 & 1 & 1.893640 & 104 & 4 & 2.970987 \\
   5 & -1 & 0.846861 & 120 & -2 & 0.691460 \\
   17 & 1 & 0.459275 & 137 & -1 & 0.161785 \\
   24 & 2 & 1.546150 & 140 & -3 & 1.440376 \\
   28 & 1 & 0.357864 & 149 & -1 & 0.155133 \\
   44 & -1 & 0.285477 & 156 & 6 & 5.458051 \\
   61 & -5 & 6.061393 & 157 & -4 & 2.418063 \\
   73 & 1 & 0.221634 & 161 & 0 & 0.000000 \\
   77 & -1 & 0.215800 & 168 & -2 & 0.584390 \\
   85 & -3 & 1.848547 & 172 & 3 & 1.299498 \\
   92 & 2 & 0.789702 & 177 & -2 & 0.569339 \\
   93 & 2 & 0.785445 & 188 & 1 & 0.138108 \\
   101 & 2 & 0.753697 & 197 & 2 & 0.539665 \\
\hline
$d$ & $c^{}_{\mf{361B}}(d)$ & $L(\mf{361B},d,1)$& $d$ & $c^{}_{\mf{361B}}(d)$ & $L(\mf{361B},d,1)$\\
\hline
   8 & 0 & 0.000000 & 97 & 0 & 0.000000 \\
   12 & -2 & 2.186587 & 105 & 2 & 0.739201 \\
   13 & 0 & 0.000000 & 109 & 0 & 0.000000 \\
   21 & -2 & 1.652904 & 113 & 2 & 0.712555 \\
   29 & 2 & 1.406560 & 124 & -8 & 10.883448 \\
   33 & -2 & 1.318562 & 129 & 2 & 0.666903 \\
   37 & 2 & 1.245250 & 136 & 4 & 2.598052 \\
   40 & 4 & 4.790572 & 141 & -2 & 0.637893 \\
   41 & -2 & 1.182947 & 145 & 2 & 0.629033 \\
   53 & 0 & 0.000000 & 165 & -2 & 0.589679 \\
   56 & -4 & 4.048772 & 173 & -2 & 0.575883 \\
   60 & -2 & 0.977871 & 181 & 6 & 5.067113 \\
   65 & 0 & 0.000000 & 184 & -4 & 2.233616 \\
   69 & 4 & 3.647479 & 185 & 2 & 0.556893 \\
   88 & -4 & 3.229803 & 193 & 4 & 2.180915 \\
   89 & 0 & 0.000000 & & &  \\
\hline
\end{tabular}
\caption{Coefficients of
$\Gross^+_{\mf{19A}}$ (top),
$\Gross^-_{\mf{19A}}$ (bottom), and central values for $\mf{361B}$
          \label{table:dlcs:19A}}
\end{table}
Table~\ref{table:dlcs:19A} shows
the values of $c^{}_{\mf{19A}}(d)$
and $L(\mf{19A},-19d,1) = L(\mf{361B},d,1)$,
where $0<d<200$ is a fundamental discriminant
such that
$19\nmid d$.
The formula
\[
  L(\mf{19A},-19d,1) = L(\mf{361B},d,1)
  = k_{\mf{19A}} \frac{c_{\mf{19A}}(d)^2}{\sqrt{d}}
\]
is satisfied, where 
\[
  k_{\mf{19A}} = \frac{2}{3}\cdot
  \frac{\<\mf{19A},\mf{19A}>}
        {L(\mf{19A},1)\sqrt{19}}
  = 1.893639859594845381072872862  = L(\mf{361B},1).
\]
% Computing M(Otilde)^f for f=361A
% sub : 10
% dps, dns : 10
% cps, cns : 0
% lps, lns : 0
% E361A x 2  : [0, 0, -1, 3, -5, 0, -7, 0]
%   e_1 = [0, 0, 0, 0, 0, 0, 0, 0, 0, 0, 1, -1, 1, -1, 1, -1, 1, -1, 1, -1, 0, 0, 0, 0, 0, 0, 0, 0, 0, 0]~  ---> 
%   e_2 = [0, 0, 0, 0, 0, 0, 0, 0, 0, 0, 0, 0, 0, 0, 0, 0, 0, 0, 0, 0, 1, -1, 1, -1, 1, -1, 1, -1, 1, -1]~  ---> 
\subsection{$\mf{361A}$}
By computing the Brandt matrices for $\Ot$,
we find the space $\M(\Ot)^{\mf{361A}}
$
of dimension $2$.
Using table~\ref{table:psi:19},
we can check that
\[
  \Gross_{\Ops}\left(\M(\Ot)^{\mf{361A}}\right)  = 
  \Gross_{\Ons}\left(\M(\Ot)^{\mf{361A}}\right)  = 0,
\]
which is expected since $L(\mf{361A},1)=0$.
% Computing M(Otilde)^f for f=361B
% sub : 10
% dps, dns : 0
% cps, cns : 0
% lps, lns : 8560
% E361B x 1  : [0, 2, 3, -1, 3, 4, -3, 0]
%   e_1 = [2, 2, 2, 2, 2, -2, -2, -2, -2, -2, -1, -1, -1, -1, -1, -1, -1, -1, -1, -1, 1, 1, 1, 1, 1, 1, 1, 1, 1, 1]~  --->  {+} 2
%   ef+     = [1, 1, 1, 1, 1, -1, -1, -1, -1, -1, -1/2, -1/2, -1/2, -1/2, -1/2, -1/2, -1/2, -1/2, -1/2, -1/2, 1/2, 1/2, 1/2, 1/2, 1/2, 1/2, 1/2, 1/2, 1/2, 1/2]~   (height = 15)
%   Theta+  = q+q^{4}+3q^{5}+3q^{9}-2q^{16}+O(q^{17})
\subsection{$\mf{361B}$}
We readily find the space $\M(\Ot)^{\mf{361B}}
=\twist\left(\M(\Ot)^{\mf{19A}}\right)$
of dimension $1$, spanned by
\begin{align*}
  \e{\mf{361B}}^+ = &\ \twist\circ\psi_{\Ot}\left(\frac{\cls{a_{1}}-\cls{a_{2}}}{2}\right),\\
\end{align*}
with height $\<\e{\mf{361B}}^+,\e{\mf{361B}}^+>=15$.
Using table~\ref{table:psi:19},
we see that
\begin{align*}
  \Gross^+_{\mf{361B}} & =  Q_{1}^{+} + 2Q_{2}^{+} + 2Q_{3}^{+} - 2Q_{4}^{+} - 2Q_{5}^{+} - Q_{6}^{+} \\ & - Q_{7}^{+} - 2Q_{8}^{+} - 2Q_{9}^{+} + 2Q_{10}^{+} + 2Q_{11}^{+} + Q_{12}^{+}; \\
\end{align*}
on the other hand
$\Gross_{\Ons}\left(\M(\Ot)^{\mf{361B}}\right)  = 0$,
which is expected since $\sign(\mf{19A})= +1$.
\begin{table}
\begin{tabular}{||r|rr||r|rr||}
\hline
$d$ & $c^{}_{\mf{19A}}(d)$ & $L(\mf{19A},d,1)$& $d$ & $c^{}_{\mf{19A}}(d)$ & $L(\mf{19A},d,1)$\\
\hline
   1 & 1 & 0.453253 & 104 & 12 & 6.400100 \\
   5 & 3 & 1.824309 & 120 & 6 & 1.489542 \\
   17 & -3 & 0.989371 & 137 & 3 & 0.348516 \\
   24 & -6 & 3.330718 & 140 & -3 & 0.344762 \\
   28 & -3 & 0.770911 & 149 & -9 & 3.007688 \\
   44 & -9 & 5.534770 & 156 & 6 & 1.306415 \\
   61 & 3 & 0.522298 & 157 & 0 & 0.000000 \\
   73 & -3 & 0.477444 & 161 & -12 & 5.143876 \\
   77 & 3 & 0.464877 & 168 & 6 & 1.258893 \\
   85 & -3 & 0.442460 & 172 & 3 & 0.311042 \\
   92 & 6 & 1.701177 & 177 & -6 & 1.226470 \\
   93 & -6 & 1.692006 & 188 & 9 & 2.677608 \\
   101 & 6 & 1.623614 & 197 & 6 & 1.162546 \\
\hline
\end{tabular}
\caption{Coefficients of $\Gross^+_{\mf{361B}}$ and central values for $\mf{19A}$
          \label{table:dlcs:361B+}}
\end{table}
Table~\ref{table:dlcs:361B+} shows
the values of $c^{}_{\mf{361B}}(d)$
and $L(\mf{361B},-19d,1) = L(\mf{19A},d,1)$,
where $0<d<200$ is a fundamental discriminant
such that
$\kro{d}{19}=+1$.
The formula
\[
  L(\mf{361B},-19d,1) = L(\mf{19A},d,1)
  = k_{\mf{361B}} \frac{c_{\mf{361B}}(d)^2}{\sqrt{d}}
\]
is satisfied, where 
\[
  k_{\mf{361B}} = \frac{19}{540}\cdot
  \frac{\<\mf{361B},\mf{361B}>}
        {L(\mf{361B},1)\sqrt{19}}
  = 0.4532532444961036035788391869  = L(\mf{19A},1).
\]
